\section{F-SM-30: the abstract interval theorem for Grothendieck categories}

\subsection{Statement and exact historical scope}
Let \(k\) be a commutative ring and let \(\mathcal M\) be a
\(k\)-linear Grothendieck abelian category with enough projectives.
Suppose that the scalar functor
\[
 \kappa:\operatorname{Mod}k\longrightarrow\mathcal M,
 \qquad U\longmapsto\mathbf1\otimes_k U,
\]
is exact, preserves coproducts, and has a left adjoint \(\varphi\).
Write \(F(X)=L\varphi(X[0])\), in nonnegative homological degrees.
Assume the following two conditions:
\begin{enumerate}
\item the canonical derived counit
\(L\varphi(\kappa k)\to k[0]\) is a quasi-isomorphism;
\item every \(X\) has a monomorphism \(X\to A_X\) with
\(L\varphi(A_X)=0\), including degree zero.
\end{enumerate}
The original full interval gives the second condition with
\(A_X=J\otimes X\).  Set
\[
 W=\{f\in\operatorname{Mor}(\mathcal M):F(f)\text{ is an equivalence}\}.
\]
Then the ordinary localization has the required equivalence
\[
 W^{-1}\mathcal M\simeq D_{\ge0}(k),                         \tag{SM30.1}
\]
induced by \(L\varphi\).  In fact the corresponding
\(\infty\)-categories are equivalent.

No exact chain-valued realization \(\Lambda\), no finite-presentation
assumption on projective generators, and no scalar projectivity
assumption on their morphism modules is required.  The additional
hypothesis relative to the bare handwritten conditions (a)--(d) is
that \(\mathcal M\) is Grothendieck: it has a generator and exact
filtered colimits.  Enough projectives is already condition (b).
The exact scalar functor and coproduct preservation follow from its
two adjoints in the original setup, as explained in F-SM-29.
We do not infer AB5 or existence of a generator from the handwriting.
Thus (SM30.1) closes this broad and precisely stated case, while
leaving a genuinely more general abelian-category assertion distinct.

\subsection{Known inputs and coherence convention}
We use the following known theorems in their stated forms.
\begin{itemize}
\item Beke (2000), Proposition 3.10, author-PDF pp.~24--25:
Grothendieck abelian categories are locally presentable.
Proposition 1.12, p.~11, gives a set generating all monomorphisms
under pushouts and transfinite composition in this setting.
Theorem 1.7 and Propositions 1.15, 1.18 provide Smith recognition
and the accessibility/solution-set implications.
\item Lurie, \emph{Higher Algebra}, Example 1.3.3.4, p.~98,
Theorem 1.3.4.4 and Proposition 1.3.4.6, p.~109:
a right-exact functor between
abelian categories with enough projectives has a genuine derived
\(\infty\)-functor, computed by projective complexes, and the
derived \(\infty\)-category is the localization at quasi-isomorphisms.
\item Lurie, \emph{Higher Algebra}, Proposition 1.3.4.22,
p.~118, and Propositions 1.3.4.24--25, p.~119: a combinatorial
model category presents a presentable \(\infty\)-category;
model homotopy colimits compute its colimits, and diagrams
indexed by ordinary categories admit strict diagram representatives.
\item Lurie, \emph{Higher Topos Theory}, Proposition 4.4.2.7,
p.~297: a functor preserving pushouts and small coproducts
preserves small colimits.  The theorem is applied for each
regular size bound, within a fixed ambient universe.
\end{itemize}
We also use Dold--Kan, Stacks Project Theorem 14.24.3,
Tag \texttt{019G}, and the projective diagram model theorem.
The normalized-total-complex calculation needed below has the
complete proof in F-SM-29, Step 3.  Its extension from modules to
a Grothendieck abelian category is explained explicitly below.
These background theorems are inputs, not claimed new proofs.

All references to \(F\) as a functor of \(\infty\)-categories
mean the actual derived functor followed by relative localization.
No functor defined only on homotopy categories is used to infer
preservation of coherent colimits.

\subsection{Complete proof}

\paragraph{1. Projectives, filtered colimits, and accessibility.}
Choose a generator \(G_0\) of \(\mathcal M\) and a projective
epimorphism \(P\to G_0\).  Then \(P\) is a projective generator:
\(\operatorname{Hom}(P,-)\) is exact, and if an object receives
no nonzero map from \(P\), it receives none from \(G_0\) and
is zero.  The evaluation map
\(\bigoplus_{\operatorname{Hom}(P,X)}P\to X\) is epic;
otherwise projectivity and generation give a map to its nonzero
cokernel which would lift and contradict evaluation.
This construction and successive kernels supply functorial
projective resolutions.  Additivity of this resolution functor
is not required.  The adjunction and exactness of \(\kappa\)
show that \(\varphi\) takes projectives to projective scalar
modules.  Coproducts of projectives are projective, since a
family of maps to an epimorphism can be lifted separately.

We verify, rather than assume, that all \(L_j\varphi\) commute
with filtered colimits.  Let \(I\) be a small filtered category.
For \(i\in I\) and projective \(Q\in\mathcal M\), the free
diagram on \(Q\) at \(i\) has value
\[
 (i_!Q)(j)=\bigoplus_{I(i,j)}Q.
\]
It is projective in \(\mathcal M^I\), by its adjunction to
exact evaluation.  These diagrams give enough projectives.
Every projective diagram is a retract of a coproduct of them.
Evaluation therefore takes projective diagrams to projectives,
and so does colimit: \(\operatorname{colim}_I i_!Q=Q\),
because the category \(i/I\) has an initial object.

Resolve an arbitrary diagram \(X:I\to\mathcal M\) by such
projective diagrams \(Q_\bullet\to X\).
Evaluation gives a projective resolution of each \(X_i\).
AB5 says colimit is exact, so \(\operatorname{colim}Q_\bullet\)
is a projective resolution of \(\operatorname{colim}X\).
The left adjoint \(\varphi\) commutes with colimits, and
filtered colimits of scalar modules are exact.  Consequently
\[
 L_j\varphi(\operatorname{colim}_I X_i)
   \simeq\operatorname{colim}_I L_j\varphi(X_i)
       \quad(j\ge0).                                    \tag{SM30.2}
\]
The isomorphism is canonical: it is induced by the diagram's
colimit maps, and comparison of projective resolutions preserves
it.  The same argument for a discrete diagram, using exact
coproducts in a Grothendieck category, proves
\[
 F\Bigl(\bigoplus_i X_i\Bigr)\simeq\bigoplus_i F(X_i).
                                                               \tag{SM30.3}
\]

The functor \(X\mapsto(L_j\varphi X)_{j\ge0}\) to graded
scalar modules is accessible, by (SM30.2) and local
presentability.  The isomorphisms in the target arrow category
form an accessible, accessibly embedded subcategory.
Beke Proposition 1.18 makes its inverse image \(W\) accessible
and accessibly embedded; Proposition 1.15 gives the solution-set
condition for any set of candidate generating cofibrations.

\paragraph{2. The model structure is constructed directly from \(L\varphi\).}
Use the set \(\mathcal I\) furnished by Beke Proposition 1.12,
so \(\operatorname{cof}(\mathcal I)\) is precisely all
monomorphisms.  Its hypotheses hold: subobjects have effective
unions in an abelian category, and AB5 preserves monomorphisms
under transfinite composition.  Closure of \(W\) under retracts
and two-out-of-three follows from its definition.

For a monomorphism \(u:A\to B\), the derived triangle of
\(0\to A\to B\to\operatorname{coker}u\to0\) gives
\[
 u\in W\quad\Longleftrightarrow\quad F(\operatorname{coker}u)=0.
                                                               \tag{SM30.4}
\]
Pushout retains this cokernel.  A transfinite composition of
monomorphisms in \(W\) stays in \(W\): use (SM30.2) at every
limit ordinal and transfinite induction on the maps from the
initial stage.  Hence the trivial cofibrations are stable under
the operations required in Smith's condition (c2).

Suppose \(u:B\to C\) has the right lifting property against
\(\mathcal I\), hence against every monomorphism.  Lifting
against \(0\to C\) gives a section, and its kernel \(K\) is
injective: a map into \(K\) extends across any monomorphism
by a lifting square with bottom map zero into \(C\).
The assumed monomorphism \(K\to A_K\) splits by injectivity.
Thus \(F(K)\) is a retract of \(F(A_K)=0\).
Since \(B\simeq C\oplus K\), the split epimorphism \(u\)
belongs to \(W\).  This is Smith's condition (c1).
All conditions now hold, so there is a combinatorial model
structure on \(\mathcal M\) with weak equivalences \(W\)
and cofibrations all monomorphisms.  Every object is cofibrant.
Let \(\mathcal C\) be its underlying \(\infty\)-category.

We use the analogous all-monomorphism, quasi-isomorphism model
on \(\operatorname{Ch}_{\ge0}(\mathcal M)\).
For clarity its existence does not result from restricting an
unbounded model.  The same recognition applies directly:
this chain category is Grothendieck, homology preserves filtered
colimits, and a monomorphism is a quasi-isomorphism exactly
when its cokernel is acyclic.  An injective chain object \(I\)
embeds into the contractible complex \(D^1(\mathbf Z)\otimes I\)
by the degree-zero inclusion; this embedding splits, so \(I\)
is acyclic.  The preceding injective-kernel argument proves
the remaining lifting condition.  This also gives the model
for \(\operatorname{Ch}_{\ge0}(k)\).
Their underlying \(\infty\)-categories are the connective
derived categories, denoted \(\mathcal D_{\ge0}(\mathcal M)\)
and \(\mathcal D_{\ge0}(k)\).

\paragraph{3. Coherent derived functors and all-scalar normalization.}
The right-exact functor \(\varphi\) gives the genuine derived
\(\infty\)-functor
\[
 \mathbb L\varphi:\mathcal D_{\ge0}(\mathcal M)
                     \longrightarrow\mathcal D_{\ge0}(k).
\]
It is the restriction of the exact right-t-exact functor of
\emph{Higher Algebra}, Example 1.3.3.4, and is computed by
applying \(\varphi\) to projective complexes.  This also
identifies its value on \(X[0]\) with the usual derived functor.
Exactness implies preservation of pushouts, which stay connective.
Projective resolutions and exact coproducts give preservation
of coproducts, with the actual comparison maps of the derived
functor.  HTT Proposition 4.4.2.7 therefore shows that
\(\mathbb L\varphi\) preserves all small connective colimits.

The exact functor \(\kappa\), preserving coproducts, induces
a functor on connective derived \(\infty\)-categories.
It too preserves pushouts and coproducts, hence all small
colimits by the same proposition.  There is the canonical
derived counit \(\mathbb L\varphi\kappa\to\mathrm{id}\).
Concretely, on a complex it is represented by
\(\varphi P(\kappa C)\to\varphi\kappa C\to C\), using
a functorial projective replacement and the original counit.
These are natural maps of complexes before localization, so
the displayed counit is a coherent natural transformation.
It is an equivalence on \(k[0]\) by hypothesis, hence on
free scalar modules and their projective retracts.
Every scalar module \(U\) is the realization of
\(\Gamma(Q_\bullet)[0]\), where \(Q_\bullet\to U\) is a
projective scalar resolution and \(\Gamma\) is Dold--Kan's
inverse normalization.  Each \(\Gamma(Q)_n\) is a finite sum
of projective modules.  Both functors in this counit preserve
that realization.  Thus the \emph{canonical} map
\[
             F(\kappa U)\longrightarrow U[0]
 \quad\hbox{is an equivalence for every }U.                \tag{SM30.5}
\]

Here and below the realization identity used is
\[
 \operatorname*{colim}_{\Delta^{\rm op}}\Gamma(C)_n[0]
                          \simeq C                       \tag{SM30.6}
\]
in the appropriate connective derived category, naturally
including its augmentation.  The proof in F-SM-29, Step 3,
works for a Grothendieck category without change to its logical
steps: normalization is an exact equivalence; coproducts and
filtered colimits are exact; the normalized chains of a simplex
are finite free abelian in each degree; their vertex contraction
acts by finite sums of identity maps on any coefficient object.
The free-diagram cell proof then gives
\(\operatorname{hocolim}Y\simeq\operatorname{Tot}NY\).
Cell cokernels, transfinite stages and retracts are handled as
there, using exact sums and AB5, and every total degree has
finitely many bidegrees.  Applying this to \(Y=\Gamma(C)[0]\)
proves (SM30.6), with its natural augmentation.

Compose \(\mathcal M\to\mathcal D_{\ge0}(\mathcal M)\),
\(X\mapsto X[0]\), with \(\mathbb L\varphi\).
By definition this coherent functor takes \(W\) to
equivalences.  The universal property of relative
\(\infty\)-localization therefore gives
\[
 \overline F:\mathcal C\longrightarrow\mathcal D_{\ge0}(k).
                                                               \tag{SM30.7}
\]
This supplies the required coherence before any colimit claim
about \(\overline F\) is made.

\paragraph{4. The localized functor preserves all colimits.}
Since every object of \(\mathcal M\) is cofibrant, its ordinary
coproducts compute model homotopy coproducts.  Equation (SM30.3)
therefore proves that \(\overline F\) preserves coproducts,
including the empty one.

Represent any pushout diagram in \(\mathcal C\) by a
cofibrant span in the model category; the strictification and
colimit comparisons are HA Propositions 1.3.4.24--25.
We may thus suppose that both maps in
\(B\leftarrow A\to C\) are monomorphisms and its ordinary
pushout is \(P=B\amalg_A C\).  There is a short exact sequence
\[
 0\longrightarrow A\xrightarrow{(u,-v)}B\oplus C
                     \longrightarrow P\longrightarrow0.
                                                               \tag{SM30.8}
\]
It exhibits the degree-zero square as a pushout in
\(\mathcal D_{\ge0}(\mathcal M)\).  Applying the actual
derived functor \(\mathbb L\varphi\) gives the pushout
square in \(\mathcal D_{\ge0}(k)\), with the same comparison
maps as (SM30.7).  Hence \(\overline F\) preserves pushouts.
Its source is presentable by HA Proposition 1.3.4.22.
HTT Proposition 4.4.2.7 now gives preservation of all small
colimits, in particular simplicial realizations.

\paragraph{5. Construction of both inverse comparisons.}
For a nonnegative scalar complex \(C\), set
\[
 L(C)=\operatorname*{colim}_{\Delta^{\rm op}}
                        \kappa\Gamma(C)_n\quad\text{in }\mathcal C.
                                                               \tag{SM30.9}
\]
Functorial projective diagram replacement may be used to choose
an ordinary representative of this colimit.  By Step 4,
(SM30.5), and (SM30.6), there is a natural equivalence
\[
                      \overline F L(C)\simeq C.            \tag{SM30.10}
\]
If \(f:C\to C'\) is a quasi-isomorphism, this comparison
shows that the chosen map representing \(L(f)\) lies in \(W\).
Thus \(L\) descends to a functor
\(\mathcal D_{\ge0}(k)\to\mathcal C\).
This argument does not assume that \(\kappa\Gamma(f)\)
is an objectwise weak equivalence.

Take the functorial projective resolution \(P_\bullet X\to X\)
from Step 1.  For a projective object \(Q\), the adjunction
unit \(\eta_Q:Q\to\kappa\varphi Q\) belongs to \(W\).
Indeed its source under \(F\) is \(\varphi Q[0]\); its
target is identified with the same object by (SM30.5); the
triangle identity says the resulting map is the identity.
In every simplicial degree \(\Gamma(P_\bullet X)_n\) is
projective.  The levelwise units consequently give equivalences
in \(\mathcal C\), and their realization gives
\[
 A_X:=\operatorname*{colim}_{\Delta^{\rm op}}
                    \Gamma(P_\bullet X)_n
       \xrightarrow{\ \sim\ }L(\varphi P_\bullet X).
                                                               \tag{SM30.11}
\]
The augmentation also gives \(a_X:A_X\to X\) in \(\mathcal C\).
We check this particular map, rather than just comparing its
source and target abstractly.  Applying \(\overline F\) and
Step 4 identifies it with the augmentation
\[
 \operatorname*{colim}_{\Delta^{\rm op}}
       \mathbb L\varphi(\Gamma(P_\bullet X)_n[0])
                   \longrightarrow\mathbb L\varphi(X[0]).
\]
Since \(\mathbb L\varphi\) preserves colimits, (SM30.6)
identifies this map with the image of the quasi-isomorphism
\(P_\bullet X\to X[0]\); hence it is an equivalence.
On ordinary representatives the augmentation map therefore
lies in \(W\), so \(a_X\) itself is an equivalence in
\(\mathcal C\).  Finally \(\varphi P_\bullet X\) is the
functorial representative of \(F(X)\).  We have natural
equivalences
\[
              X\xleftarrow{\ \sim\ } A_X
                    \xrightarrow{\ \sim\ }L(\overline F X).
                                                               \tag{SM30.12}
\]
The augmentations, units and resolution maps used to construct
this span are natural.  Equations (SM30.10)--(SM30.12) prove
that \(L\) and \(\overline F\) are mutually inverse
\(\infty\)-functors.  Passing to homotopy categories gives
(SM30.1). \(\square\)

\subsection{What the proof removes and what it does not remove}
The proof constructs accessibility and the model structure
directly from \(L\varphi\), rather than adding an exact
chain-valued realization as an assumption.  It works for every
commutative scalar ring and for projective generators which
need not be finitely presented.  For the underlying localization,
the full tensor interval is used only through the stated
acyclic-monomorphism consequence.

The AB5 step is substantive.  Merely observing that an
endofunctor such as \(V\mapsto V^{**}/V\) is exact over a field
does not make its square-zero module category Grothendieck.
For the intake's three-sort proposal, an asserted periodic
formula \(L_{3n}\varphi(\kappa V)=(V^{**}/V)^{\circ n}\)
would contradict (SM30.5) for infinite-dimensional \(V\),
so all of the present hypotheses cannot simultaneously hold
there.  This is a diagnostic conditional implication, not a
completed reconstruction of that counterexample: its symmetric
tensor coherence and full interval still require their absent
proof.  No conclusion about that missing proof is inferred.

There is also no ordinary right adjoint \(G\) claimed here.
Whenever one chooses an exact chain-valued adjunction separately,
its raw counit remains the stronger and separately identified
test in BR-13403.  The present theorem concerns the underlying
localization in the explicit Grothendieck range.

\subsection{Primary references}
\begin{itemize}
\item Beke, T. (2000). Sheafifiable homotopy model categories.
\emph{Mathematical Proceedings of the Cambridge Philosophical
Society, 129}(3), 447--475.
\url{https://doi.org/10.1017/S0305004100004722}.
Author PDF: \url{https://faculty.uml.edu/tbeke/sheafi.pdf}.
\item Lurie, J. (2009). \emph{Higher topos theory}
(Annals of Mathematics Studies, Vol.~170). Princeton University
Press. \url{https://www.math.ias.edu/~lurie/papers/HTT.pdf}.
Exact input: Proposition 4.4.2.7, printed p.~297.
\item Lurie, J. (2017, September 18). \emph{Higher algebra}.
Author manuscript.
\url{https://people.math.harvard.edu/~lurie/papers/HA.pdf}.
Exact inputs: Example 1.3.3.4, p.~98; Theorem 1.3.4.4 and
Proposition 1.3.4.6, p.~109; Proposition 1.3.4.22, p.~118; and Propositions
1.3.4.24--25, p.~119.
\item The Stacks Project Authors. (2026). \emph{The Stacks
Project}. Theorem 14.24.3, Tag \texttt{019G}.
\url{https://stacks.math.columbia.edu/tag/019G}.
\end{itemize}
